\documentclass[11pt]{article}
\usepackage[margin=0.85in]{geometry}
\usepackage{enumitem}
\usepackage{xcolor}
\usepackage{amsmath}
\usepackage{booktabs}
\usepackage{array}
\usepackage{longtable}

% Colors (match diagram)
\definecolor{episodecolor}{RGB}{70, 70, 70}
\definecolor{intervalcolor}{RGB}{31, 119, 180}
\definecolor{roundcolor}{RGB}{44, 160, 44}
\definecolor{elementcolor}{RGB}{214, 39, 40}
\definecolor{solvercolor}{RGB}{148, 103, 189}
\definecolor{calloutcolor}{RGB}{255, 127, 14}
\definecolor{codeblue}{RGB}{0, 100, 180}

\newcommand{\method}[1]{\textcolor{codeblue}{\texttt{#1}}}

\pagestyle{empty}

\begin{document}

\begin{center}
{\Large\bfseries Code Reference: Architecture $\to$ Implementation}
\end{center}

\vspace{0.15cm}
\noindent\textit{Companion to the architecture description artifacts. Maps each architectural concept to the method(s) that implement it. Organized by component so you can flip directly to the relevant section during discussion. For full implementation detail, see the source code.}

\vspace{0.3cm}

% ============================================================
% KEY DATA STRUCTURE
% ============================================================
\noindent\textbf{Key Data Structure: \method{solver.label\_mat}}

\vspace{0.15cm}

\noindent The binary tree that encodes the mesh hierarchy. Every element that has ever existed (active or not) has a row in this matrix. Each row contains:

\begin{center}
\renewcommand{\arraystretch}{1.3}
\begin{tabular}{c l l}
\toprule
Column & Field & Meaning \\
\midrule
0 & \texttt{self\_id} & This element's unique ID (1-based) \\
1 & \texttt{parent\_id} & Parent element's ID (0 if base level) \\
2 & \texttt{child1\_id} & Left child's ID (0 if leaf / never refined) \\
3 & \texttt{child2\_id} & Right child's ID (0 if leaf / never refined) \\
4 & \texttt{level} & Refinement level (0 = base mesh) \\
\bottomrule
\end{tabular}
\end{center}

\noindent Only elements in \method{solver.active} are currently part of the mesh. An element's children replace it in \method{active} when it is refined; the parent replaces both children when they are coarsened. The tree structure is persistent --- refined-then-coarsened elements retain their entries, enabling re-refinement without rebuilding.

\vspace{0.1cm}
\noindent \method{solver.active} is an array of element IDs (1-based) in left-to-right spatial order.

% ============================================================
% COMPONENT 3: ELEMENT VISIT
% ============================================================
\vspace{0.5cm}
\begin{center}
\rule{0.6\textwidth}{0.4pt}
\end{center}
\vspace{0.3cm}

\noindent{\large\textbf{\textcolor{elementcolor}{Component 3: The Element Visit}}}

\vspace{0.3cm}

% --- action_masks ---
\noindent\textbf{\method{action\_masks()} \textnormal{--- MaskablePPO interface}}

\vspace{0.1cm}
\noindent Returns a 3-element boolean array indicating which actions (coarsen, do-nothing, refine) are valid for the current element. Called by MaskablePPO before each action selection. Do-nothing is always valid. Refine is blocked only when the element is at \texttt{max\_level}. Coarsen validity is delegated to \method{\_can\_coarsen()}.

\vspace{0.3cm}

% --- _can_coarsen ---
\noindent\textbf{\method{\_can\_coarsen(active\_idx)} \textnormal{--- coarsen validation (3 conditions)}}

\vspace{0.1cm}
\noindent Checks the three structural conditions for a valid coarsen action:

\begin{enumerate}[leftmargin=1.5em, itemsep=3pt]
    \item \textbf{Sibling exists and is active.} Delegates to \method{\_find\_sibling()} (see below). Returns \texttt{False} if sibling is \texttt{None}.
    \item \textbf{Sibling is not cascade-created this round.} Checks whether the sibling's element ID is in \method{self.consumed\_elements}, the set of IDs created by balance cascades during the current round. Prevents the agent from undoing balance enforcement.
    \item \textbf{Post-coarsening 2:1 balance.} Computes what the parent's level would be after coarsening (\texttt{current\_level $-$ 1}), then checks that both external neighbors (left of the left sibling, right of the right sibling, with periodic wrap) have levels within 1 of the parent level. Uses \method{\_get\_element\_level()} for each neighbor.
\end{enumerate}

\vspace{0.3cm}

% --- _find_sibling ---
\noindent\textbf{\method{\_find\_sibling(active\_idx)} \textnormal{--- binary tree sibling lookup}}

\vspace{0.1cm}
\noindent Finds the sibling of an element using the binary tree:
\begin{enumerate}[leftmargin=1.5em, itemsep=2pt]
    \item Looks up the element's \texttt{parent\_id} from \method{label\_mat}.
    \item If \texttt{parent\_id} is 0, the element is at base level --- return \texttt{None}.
    \item Gets both children of the parent (\texttt{child1\_id}, \texttt{child2\_id}) from the parent's row in \method{label\_mat}.
    \item The child that isn't the current element is the sibling.
    \item Searches \method{solver.active} for the sibling's ID. If found, returns its index; if not (sibling has been further refined and replaced by its own children), returns \texttt{None}.
\end{enumerate}

\vspace{0.3cm}

% --- _get_element_level ---
\noindent\textbf{\method{\_get\_element\_level(active\_idx)} \textnormal{--- refinement level query}}

\vspace{0.1cm}
\noindent Reads column 4 of the element's row in \method{label\_mat}. One line: looks up the element ID from \method{solver.active[active\_idx]}, indexes into \method{label\_mat}, returns the integer level.

\vspace{0.3cm}

% --- _execute_action ---
\noindent\textbf{\method{\_execute\_action(active\_idx, action)} \textnormal{--- action execution + cascade tracking}}

\vspace{0.1cm}
\noindent Executes the agent's chosen action in two deliberate stages:

\begin{enumerate}[leftmargin=1.5em, itemsep=3pt]
    \item \textbf{Apply the mesh action without balance.} Calls \method{solver.adapt\_mesh()} with \texttt{balance=False}. For refinement, the solver replaces the element with two children using prolongation matrices. For coarsening, both siblings are merged into their parent using restriction matrices.
    \item \textbf{Enforce balance separately.} Calls \method{solver.balance\_mesh()}. Any additional refinements triggered by balance violations happen here.
    \item \textbf{Detect cascades.} Calls \method{\_detect\_cascade\_elements()}, which compares the set of active element IDs before and after balance enforcement. The set difference identifies cascade-created elements. These IDs are added to \method{self.consumed\_elements} for the remainder of the round.
\end{enumerate}

\noindent The two-stage separation is critical: it lets the environment distinguish the agent's decision from forced balance corrections, which matters for reward attribution, queue management, and action masking.

\vspace{0.3cm}

% --- _detect_cascade_elements ---
\noindent\textbf{\method{\_detect\_cascade\_elements(pre, post)} \textnormal{--- set-difference cascade detection}}

\vspace{0.1cm}
\noindent Takes two sets of element IDs: the active set after the agent's action (before balance) and the active set after balance enforcement. Returns \texttt{post $-$ pre} --- any element IDs that appeared as a result of balance enforcement are cascade-created elements.

\vspace{0.3cm}

% --- _advance_queue ---
\noindent\textbf{\method{\_advance\_queue()} \textnormal{--- queue traversal and transition handling}}

\vspace{0.1cm}
\noindent Moves to the next element in the queue. The queue stores element IDs (not indices) for stability across mesh changes. At each advance:
\begin{enumerate}[leftmargin=1.5em, itemsep=2pt]
    \item Increments \method{queue\_position}.
    \item Resolves the next element ID to a current \method{solver.active} index. If the ID is no longer in the active list (element was consumed by a cascade or coarsened away), skips it and tries the next.
    \item If the queue is exhausted and more rounds remain, increments \method{round\_number}, clears \method{consumed\_elements}, rebuilds the queue via \method{\_build\_queue()}, and continues.
    \item If all rounds are exhausted, returns a transition signal (\texttt{`interval'} or \texttt{`done'}) to the caller (\method{step()}), which then handles the solver advance and global reward.
\end{enumerate}

% ============================================================
% COMPONENT 4: OBSERVATION SPACE (placeholder)
% ============================================================
\vspace{0.5cm}
\begin{center}
\rule{0.6\textwidth}{0.4pt}
\end{center}
\vspace{0.3cm}

\noindent{\large\textbf{\textcolor{elementcolor}{Component 4: Observation Space}}}

\vspace{0.15cm}
\noindent\textit{To be added when Component 4 artifact is developed.}

% ============================================================
% COMPONENT 5: QUEUE MECHANICS (placeholder)
% ============================================================
\vspace{0.5cm}
\begin{center}
\rule{0.6\textwidth}{0.4pt}
\end{center}
\vspace{0.3cm}

\noindent{\large\textbf{\textcolor{roundcolor}{Component 5: Queue Mechanics}}}

\vspace{0.15cm}
\noindent\textit{To be added when Component 5 artifact is developed.}

% ============================================================
% COMPONENT 6: DUAL REWARD STRUCTURE (placeholder)
% ============================================================
\vspace{0.5cm}
\begin{center}
\rule{0.6\textwidth}{0.4pt}
\end{center}
\vspace{0.3cm}

\noindent{\large\textbf{\textcolor{solvercolor}{Component 6: Dual Reward Structure}}}

\vspace{0.15cm}
\noindent\textit{To be added when Component 6 artifact is developed.}

\end{document}
